% =====================================================================
%  PART VI -- THE VERDICT, THE DISCLOSURE, AND THE APPENDIX
%  (merged edition v3, section 6)
% =====================================================================
\clearpage
\part{The verdict, the disclosure, and the appendix}

\section{The verdict}
\label{sec:verdict}

\subsection{Clay: \texorpdfstring{$0\%$}{0 percent}}
\label{sec:verdict-clay}

This subsection states the figure first and the reasons second, because the figure is the
deliverable and the reasons are only its justification.

\begin{claimbox}{cN}
\textbf{Clay: $0\%$.} The four-dimensional pure $SU(N)$ Yang--Mills mass gap is not proved in
this paper, is not reduced to a finite computation in this paper, and is not placed within one
lemma of any tool developed here. What is proved is a reduction, a set of supporting theorems, one
measurement at the physical coupling, and a systematic negative survey. The percentage above is a
judgement about distance to a complete unconditional proof, not a measure of effort.
\end{claimbox}

The honest structure of the failure is not ``we did not find the argument''. It is that
\textbf{each of the three legs of the reduction closes on its own premise}: the reduction is exact,
and every attempt to discharge a leg returns that leg. We take the three legs in turn.

\subsubsection*{(i) Pointwise positivity at the physical coupling}

The leg asks for $w_1(2.5) > 0$ for the physical operator. \textbf{Numerically, it holds.} On
$L_t = 12$, with the \code{njit} heat bath at $\beta = 2.5$, on $L_s = 6$ and $L_s = 8$, the
measurement gives
\[
  w_1(2.5) \;=\; 0.3427 \,\pm\, 0.0119 \qquad (28.8\,\sigma \text{ from zero}),
\]
so the sign is not in doubt at this lattice spacing and this volume. But the statement is not a
statement about the Hamiltonian: on the \emph{same} Markov chain and the same $12000$
configurations, a second operator gives $0.0168 \pm 0.0142$. \textbf{The two differ by a factor
$20.2$.} The one-line counterexample $w_1 = \cos^2\theta$ at fixed spectrum shows why: $w_1$ is a
property of the pair $(H, O)$, not of $H$. Hence the leg splits in two: its pointwise form is \LE{}
(numerically satisfied), while its uniform form is \LO{} and, until $O$ is named, vacuous as a
statement about the theory. Transporting it to the continuum requires an $O$-independence statement
that does not exist here.

\subsubsection*{(ii) The $L$-monotone step}

The leg asks for a monotone quantity in the volume that survives the infinite-volume limit.
\textbf{No such step exists in the class of quantities considered.} The obstruction is two-sided
realizability: both a $2/(L-1)$ profile and a constant profile are realizable within the same
structural assumptions, so no monotonicity theorem can be extracted from the structure alone; and
there is \textbf{no FKG inequality} for the lattice gauge measure at $\beta = 2.5$ to supply the
comparison that realizability leaves open. The leg is therefore not merely unproved but
unavailable: nothing in the present toolkit produces it. \LN{}

\subsubsection*{(iii) The $\beta$-tail-uniform floor}

The leg asks for a floor on the spectral quantity that is uniform in $\beta$ along the tail.
\textbf{This floor is $\Delta_0$, and $\Delta_0$ has already been judged dead}: the model $M6$ is
insufficient to supply it, and the two-dimensional theory is unnecessary for it, so neither the
solvable-model route nor the dimensional-reduction route reaches it. Renaming the floor does not
create one; the leg and the target are the same object. \LO{}

\begin{v3box}
\textbf{Reading of the three legs together.} (i) holds numerically and depends on $O$ with a factor
$20$; (ii) does not exist as a step, and no FKG is available to create it; (iii) is $\Delta_0$,
already dead. The three failures are of three different kinds --- empirical-but-operator-bound,
structurally-absent, and identically-the-target --- which is why the reduction does not collapse
into a single missing lemma.
\end{v3box}

\src{Leg (i): the measurement of \S\ref{sec:measurement}, $\beta = 2.5$, $L_t = 12$; operator
pair as stated in \S\ref{sec:m-opdep}. Leg (ii): two-sided realizability construction,
\code{research/unconditional/} notes. Leg (iii): $\Delta_0$ and its ledger entry,
\S\ref{sec:ledger}.}

\subsubsection*{The direction of the twenty routes}

The verdict is reinforced, rather than merely asserted, by the survey of \S\ref{sec:survey}: twenty
routes to uniformity in $(L, a)$ were attacked, and they fail in a \emph{uniform direction}. That
uniformity is the substantive finding, not the individual failures.

\begin{center}
\small
\begin{tabularx}{\textwidth}{@{}l X c@{}}
\toprule
\textbf{Class} & \textbf{Routes} & \textbf{Count} \\
\midrule
Upper bound only & compactness, reflection positivity, Berezin--Lieb, the probabilistic reading,
the Krein--Rutman residual, concavity, the RG modulus, two Ursell constants & $9$ \\
Circular & positivity improvement, minimality, operator inequalities, definability, the
strong--weak bridge, and the named obstructions $\Delta_0$, $T4$, $H1$, $M(\beta)$ & $9$ \\
Counterexample & two natural implications, each killed by an explicit machine-checked
counterexample & $2$ \\
\bottomrule
\end{tabularx}
\end{center}

\noindent The direction is one-way: \textbf{the upper side is complete and the lower side is
empty.} Nine routes give only upper bounds on the quantity wanted; nine give lower bounds only by
assuming what is to be proved; two are false as stated. Two routes did produce something usable ---
a block-Loewner inequality that repairs a previously unproved step, and a sampler bug whose
correction invalidates four earlier numerical sign determinations --- but neither is a lower bound.
A survey in which every independent route fails in the same direction is evidence about the shape
of the problem, not a list of unlucky attempts.

\subsubsection*{The only framework that survives, and what it gives}

Exactly one piece of general machinery survives the survey without circularity: the pair
\textbf{Banach--Alaoglu} (for the existence of weak limits of the finite-volume measures) together
with \textbf{Portmanteau} (for the identification of the limit's expectation functionals). Both
are free, and both are free in the direction that does not help.

\begin{itemize}
\item \textbf{Tightness is free.} The lattice measures are tight by construction on the compact
field space, so the compactness class of arguments is excluded as a source of new information: it
supplies existence of a limit, never a gap.
\item \textbf{Upper semicontinuity is free.} Portmanteau then gives the inequality that runs from
finite volume to the limit, and only in that direction. There is no companion inequality running
from the limit back to finite volume, and the gap is exactly the missing companion.
\end{itemize}

\noindent Consequently the free machinery yields precisely the direction that the nine upper-bound
routes already used, and says nothing about the lower side. This is the technical content of the
sentence ``the wall is a direction, not a difficulty''.

\subsection{What would be needed}
\label{sec:verdict-needed}

The reduction of \S\ref{sec:reduction-equiv} identifies the missing object, and it is worth naming
its \emph{type} rather than its content, because every previous attempt in this project named the
content and thereby restated the target.

\begin{qbox}
\textbf{What is missing is one machine that produces lower bounds} --- not another transport of an
upper bound across the same interface. The distinction is not rhetorical: the nine upper-bound
routes of \S\ref{sec:survey} all share a common form (a free inequality applied to a free
compactness or semicontinuity statement), and the nine circular routes all share the complementary
form (a lower bound obtained by assuming a lower bound). A machine of the missing type must produce
a lower bound from hypotheses that do not include one.
\end{qbox}

Every attempt in this project to build that machine terminated at the same point, and it terminated
there \emph{independently} in three places:

\begin{enumerate}
\item \code{N09-1}: the target is unique --- any object strong enough to be transported is already
the gap, so there is no strictly weaker surrogate to transport.
\item \code{u04}: the non-circular version of the argument \emph{is} the Clay problem --- removing
the circular step removes the argument, not the circle.
\item \code{mon1}: the point $L = \infty$ is itself the gap --- the limit object whose positivity
would be transported is the object whose positivity is wanted.
\end{enumerate}

\noindent Three independent lines of work, arriving at three different formulations, identify the
same point. We regard that coincidence as the strongest single piece of structural evidence
produced by the campaign, and it is the reason the verdict is stated as a direction rather than as
a shortage of effort.

\begin{qbox}
\textbf{This subsection is not a restatement of \S\ref{sec:verdict-clay}, and the distinction is
deliberate.} \S\ref{sec:verdict-clay} lists \emph{which} legs are closed and \emph{why} each closure
is of its own kind; this subsection states \emph{what type of object} is absent, and records the
three independent identifications of that object with the target. The first is a status report on
three named legs; the second is a claim about the shape of any future proof. Where the two would
have overlapped --- the naming of $\Delta_0$ as the third leg and as one of the nine circular routes
--- the material is stated once, in \S\ref{sec:verdict-clay}, and referenced here.
\end{qbox}

\subsection{What is nevertheless new}
\label{sec:verdict-new}

The negative verdict does not license self-effacement. Five results below are new, are owned by
this paper, and would survive the disappearance of every other claim in it.

\begin{enumerate}
\item[\textbf{1.}] \textbf{A window correction of a factor $4.87$.} A previously published
criterion window, $0.018394$, is the $a = 0$ activation root and is illegal as a criterion window;
the correct value is $3.777\times 10^{-3}$, and $0.018394/0.003777 = 4.87$ (the companion ratio
$4.84$ belongs to the earlier $0.018269$ of \S\ref{sec:est-window-fix}). The over-report was found
and corrected by this project
against its own published claim. A published over-report that is caught by its own author is a
result, and we do not downgrade it to an erratum.
\item[\textbf{2.}] \textbf{The first numerical value of $w_1$ at $\beta = 2.5$.} The measurement
$w_1 = 0.3427 \pm 0.0119$ at $28.8\sigma$, on $L_s = 6$ and $L_s = 8$ at $L_t = 12$, with the
operator dependence ($0.0168 \pm 0.0142$ for the second operator, factor $20.2$) explicitly
labelled rather than suppressed. The labelling is part of the result: an unlabelled version of this
number would be misleading.
\item[\textbf{3.}] \textbf{The no-go of \code{SG1-1}} (\LC{}). A sixth species of proof is
excluded, at grade \LC{} rather than \LT{} because one step does not survive a computability
reading. The downgrade from \LT{} to \LC{} was reported by the author of that result, not by a
critic.
\item[\textbf{4.}] \textbf{A systematic negative survey of twenty routes} (\S\ref{sec:survey}),
organized by the direction of failure rather than by chronology. Its value is the invariant
direction --- nine upper, nine circular, two counterexamples --- not the count.
\item[\textbf{5.}] \textbf{A ledger of thirty-eight definitional collisions} (\S\ref{sec:ledger}).
This is a methodological contribution: it is the mechanism by which the errors of this paper were
caught, including twice during the writing of this manuscript.
\end{enumerate}

\noindent Item 5 is not bookkeeping decoration. Two of the thirty-eight entries --- corpus item~35,
a negative $r_0$ (row~26 of Table~\ref{tab:ledger}), and corpus item~36, an over-broad formulation
of the asymptotic-perturbation estimate (row~27) --- were found
\emph{while this section was being drafted}, and both were found by an independent review pass
rather than by the author of the affected claim. That is the mechanism, and it is the subject of
the next subsection.

\subsection{AI contribution statement}
\label{sec:verdict-ai}

This paper is a multi-agent artifact, and the statement below is written so that a reader can tell
what was done by which kind of agent.

\begin{claimbox}{cC}
\textbf{The great majority of the derivations, the numerical work, and the audits in this paper were
carried out by AI agents}, over many rounds and across more than a hundred sub-agent assignments.
\textbf{The human author's role was fourfold}: to set the target, to hold the honesty line, to
refuse to call an unproved statement a theorem, and to demand, in every round, an answer to the
question \emph{who is opposing this}. The last of these is the one that mattered most, because it
converted each round from a claim into a contest.
\end{claimbox}

Concretely:

\begin{itemize}
\item \textbf{All thirty-eight ledger entries were caught by independent review} --- AI auditing
AI, with the reviewer separated from the author of the claim. No entry was caught by the agent that
made the error.
\item \textbf{Two of them were caught during the writing of this manuscript}: corpus items 35 and 36
(rows 26 and 27 of Table~\ref{tab:ledger}) --- a negative $r_0$ and an over-broad
asymptotic-perturbation estimate. Both were introduced in earlier
rounds and survived until this section was drafted.
\item \textbf{The downgrade of \code{SG1-1} from \LT{} to \LC{}} was reported by its own author,
after a computability reading broke one step of the argument.
\end{itemize}

\begin{warnbox}
\textbf{Framing, stated explicitly.} The sentence that carries this subsection is not ``AI agents
produced correct mathematics'' and must not be read as a claim of credit for the result. It is:
\emph{the reliability of this paper does not come from AI agents being unable to err; it comes from
the requirement that every error be caught by an independent path.} The thirty-eight ledger entries
are evidence for the second sentence and evidence against the first: an error process that
produced thirty-eight caught collisions is not an error process that can be trusted to have
produced no uncaught ones. The human author's contribution is therefore best described as the
maintenance of that requirement, and the paper's claim to reliability is a claim about the
requirement, not about the agents.
\end{warnbox}

\noindent We also record the negative side of the disclosure: no step of the reduction, no window
correction, and no grade assignment in this paper was accepted on the authority of a single agent,
and the places where a single agent's conclusion was allowed to stand are exactly the places marked
\LE{} --- measured, not proved.

\appendix

\section{Reproducibility}
\label{sec:appendix-repro}

\noindent\emph{This appendix is the reproducibility section of the merged draft; it prints as
Appendix~A because the source issues \code{\textbackslash{}appendix} here.}

Everything numerical, formal, and enumerative in this paper is regenerable from the artefacts
listed here. Paths of all scripts and logs are given in their respective \code{*.md} files, which
are the authoritative index; the summary below is for orientation only.

\subsection*{Numerical measurement}

\begin{itemize}
\item \textbf{Algorithm}: \code{njit} heat bath, single long chain, thermalized once, no
re-thermalization between volumes.
\item \textbf{Parameters}: $\beta = 2.5$, $L_s = 6$ and $L_s = 8$, $L_t = 12$, $12000$
configurations on the shared chain.
\item \textbf{Cost}: $2.223$~ms/sweep at $L = 6$ and $5.390$~ms/sweep at $L = 8$, measured on the
reference machine; the long-chain design means the $L = 8$ run is the binding cost.
\item \textbf{Outputs}: $w_1 = 0.3427 \pm 0.0119$ and $0.0168 \pm 0.0142$ for the two operators
(\S\ref{sec:measurement}); the sampler bug that invalidates four earlier sign determinations is documented with its
fix in the same directory.
\end{itemize}

\subsection*{Formal verification}

\begin{itemize}
\item \textbf{File}: \code{M02-CBAR.lean}.
\item \textbf{Content}: $25$ declarations, \code{sorry = 0}.
\item \textbf{Axioms}: \code{[propext, Classical.choice, Quot.sound]} --- no axiom beyond the three
permitted by the grade rules stated on the title page.
\item \textbf{Scope}: the isolation statement and its supporting lemmas; the reduction itself is
not formalised.
\end{itemize}

\subsection*{Exhaustive enumeration}

\begin{itemize}
\item $a_7 = 189\,215\,964$, exact, by complete enumeration rather than by extrapolation.
\item $b(\Gamma) \le 14n + 6$, and the constant is sharp.
\item $t_n \le B^n$ with $B = 19^{19}/18^{18} = 50.2820906$.
\end{itemize}

\subsection*{Provenance}

Every number quoted above appears with its source file at the point of use in
\S\ref{sec:established}--\S\ref{sec:ledger}. Where a
number is a judgement rather than a measurement --- the $0\%$ of \S\ref{sec:verdict-clay}, the $20$
in ``twenty
routes'', the count of thirty-eight ledger entries --- it is labelled as such at the point of use.
The provenance rule stated on the title page is the one rule this paper does not relax in its own
favour.
